% Generated draft: requires mathematical and source review.
\documentclass[11pt]{article}
\usepackage[margin=25mm]{geometry}
\usepackage{fontspec}
\setmainfont{Latin Modern Roman}
\usepackage{amsmath,amssymb,mathtools,unicode-math}
\setmathfont{Latin Modern Math}
\usepackage{enumitem,needspace,xurl,hyperref}
\setlength{\parindent}{0pt}
\setlength{\parskip}{5pt}
\setlength{\emergencystretch}{3em}
\newcommand{\N}{\mathbb{N}}
\newcommand{\Z}{\mathbb{Z}}
\newcommand{\Q}{\mathbb{Q}}
\newcommand{\R}{\mathbb{R}}
\newcommand{\C}{\mathbb{C}}
\newcommand{\F}{\mathbb{F}}
\newcommand{\A}{\mathbb{A}}
\newcommand{\PP}{\mathbb{P}}
\newcommand{\E}{\mathbb{E}}
\newcommand{\eps}{\varepsilon}
\newcommand{\dd}{\,\mathrm{d}}
\begin{document}
\section*{UnsolvedMath ID: 30005017}
\textbf{Source problem code:} catalogue entry
\subsection{Definitions and mathematical statement}
Let \(d\geq 1\), and let \(P\subset\R^d\) be a polytope with vertices \(v_1,\ldots,v_n\). The principal question in the supplied record concerns \(d=3\). Write \(\lambda_1,\ldots,\lambda_n\) for the Wachspress coordinate functions associated with these vertices, wherever the chosen construction defines them. The source treats these functions as rational functions. Their barycentric normalization, wherever defined, is
\[
\sum_{i=1}^n\lambda_i(t)=1,
\qquad
\sum_{i=1}^n\lambda_i(t)v_i=t.
\]
These identities describe the barycentric normalization; they do not uniquely define the Wachspress functions. The supplied record gives neither their construction nor additional hypotheses on the polytope under which that construction is intended.

Here is a precise formulation conditional on a specified rational Wachspress construction. Let \(K\) be either \(\R\) or \(\C\), using the same real rational functions over \(\C\) in the latter case. Let \(U_K\subset\A_K^d\) be a nonempty Zariski open subset on which all the functions \(\lambda_i\) are regular and do not vanish simultaneously. They define the rational map
\[
\Phi_{P,K}\colon\A_K^d\dashrightarrow\PP_K^{n-1},
\qquad
t\longmapsto[\lambda_1(t):\cdots:\lambda_n(t)].
\]
Set
\[
X_{P,K}=\overline{\Phi_{P,K}(U_K)}^{\,\mathrm{Zar}}
\subset\PP_K^{n-1}.
\]
With homogeneous coordinates \(z_1,\ldots,z_n\) indexed by the vertices, its homogeneous vanishing ideal is
\[
I_{P,K}=\bigoplus_{m\geq0}
\left\{F\in K[z_1,\ldots,z_n]_m:
F(\lambda_1(t),\ldots,\lambda_n(t))=0
\text{ in }K(t_1,\ldots,t_d)\right\}.
\]
\textbf{Problem.} For the intended class of three-dimensional polytopes and its Wachspress construction, determine homogeneous generators of \(I_{P,K}\) in terms of \(P\), and describe the associated projective variety \(X_{P,K}\). Formulate and determine the corresponding equations and geometric descriptions in higher dimensions \(d>3\).

The source calls the variety three-dimensional and states that its degree is known, without supplying a degree formula or its hypotheses. Full dimensionality of \(P\) is needed for the stated three-dimensional interpretation but is not explicit in the supplied hypothesis \(P\subset\R^3\). The source also does not select the real or complex ground field, or specify restrictions such as simplicity of \(P\). The formulation above therefore keeps the ground field explicit and is conditional on a well-defined rational coordinate construction. It uses the rational-map domain and does not prescribe values at boundary points or poles where the coordinate expressions are undefined.

An opening fragment of the supplied original statement precedes the heading about Wachspress varieties. It mentions spans corresponding to the five partitions of \(4\):
\[
\begin{gathered}
K\{1,x,x^2,x^3\},\qquad K\{1,x,x^2,y\},\qquad
K\{1,x,y,xy\},\\
K\{1,x,y,y^2\},\qquad K\{1,y,y^2,y^3\}.
\end{gathered}
\]
Here the displayed notation denotes linear spans of the listed monomials over an unspecified field \(K\). The fragment says that these provide an upper bound, that at least three are needed, and that four is conjectured to be the correct number. It does not identify the object being bounded or establish a connection to the Wachspress question. These numerical assertions cannot be made part of the mathematical target above without further source context.
\subsection{Short English statement}
Given a three-dimensional polytope and a specified rational Wachspress coordinate construction, find homogeneous polynomial equations defining the projective Zariski closure of its coordinate image, describe that variety, and seek analogous descriptions in higher dimensions. The supplied record leaves the precise polytope hypotheses and ground field unresolved and includes an apparently separate fragment whose intended target is unspecified.
\subsection{Sources}
\begin{itemize}
\item Supplied local UnsolvedMath registry record, permanent ID \(30005017\), exact source problem code \texttt{catalogue entry}, title: \emph{Equations of Three-Dimensional Wachspress Varieties}. The record labels the problem open and its statement \texttt{corrected\_verified}; its claim of checking against the official report is supplied provenance, not independent verification here.
\item Supplied source citation: \emph{Mini-Workshop: Interpolation, Approximation, and Algebra (2023)}, Oberwolfach Reports, pp.\ 383--428, DOI \texttt{10.4171/owr/2022/7}. \url{https://doi.org/10.4171/owr/2022/7}. The original statement attributes the three-dimensional Wachspress-variety questions to Sottile and Schenck and reports that Schenck, Smith, and Sottile investigated them in 2013 without a conclusive result. The registry separately records proposed year 2022.
\item Supplied literature source: OWR barycentric-coordinate problem session. \url{https://doi.org/10.4171/owr/2022/7}.
\item Supplied literature source: Irving--Schenck, related Wachspress surfaces. \url{https://doi.org/10.2140/ant.2014.8.369}.
\item Literature assessment dated exactly \texttt{2026-08-22}: supplied dated curation, not independently verified. It labels the problem open and says that algebraicity and basic properties are known for mean-value barycentric coordinates, while a general set of homogeneous equations for the requested Zariski closure remains open. This assessment concerns mean-value coordinates, whereas the supplied problem concerns Wachspress coordinates; its applicability is unresolved. The background explicitly requests independent verification of time-sensitive status and cited claims.
\end{itemize}
\end{document}
